\documentclass[12pt, a4paper]{article}
\usepackage[utf8]{inputenc}
\usepackage[T1]{fontenc}
\usepackage[spanish]{babel}
\usepackage[left=3.5cm, right=3.5cm, top=2.5cm, bottom=2.5cm]{geometry}
\usepackage{parskip}
\usepackage{titlesec}
\usepackage{hyperref}
\usepackage{lmodern}
\usepackage{booktabs}

\titleformat{\section}
  {\large\bfseries}{}{0em}{}[\vspace{-0.4em}\rule{\linewidth}{0.6pt}\vspace{0.2em}]

\hypersetup{
    colorlinks=true,
    linkcolor=black,
    urlcolor=black,
    pdftitle={Autoevaluacion Entregable 4 - GEP},
}

\setlength{\parskip}{0.5em}

% ------------------------------------------------------------------
\begin{document}

% ---- Portada ----
\begin{titlepage}
    \centering
    \vspace*{3cm}
    {\LARGE\bfseries Autoevaluación del Entregable 4\\[0.5em]}
    {\large Gestió de Projectes\\[3em]}
    {\normalsize
        \textbf{Autor:} Adrià Cebrián Ruiz\\[0.4em]
        \textbf{Director:} Luis Jonás Cuervo Derqui\\[0.4em]
        \textbf{Ponente:} David Garcia Soriano\\[0.4em]
        \textbf{Tutor GEP:} Joan Sardà Ferrer\\[2em]
        Facultat d'Informàtica de Barcelona (FIB)\\
        Universitat Politècnica de Catalunya\\
        Menció de Computació · Grau en Enginyeria Informàtica\\[2em]
        20 de marzo de 2026
    }
\end{titlepage}

% ---- Introducción ----
\noindent Este documento recoge la autoevaluación del Entregable 4, el documento final integrado de GEP. Este entregable consolida y corrige los tres entregables anteriores, incorporando el \textit{feedback} recibido del profesor a lo largo del curso.

% ==================================================================
\section{Contextualització i Abast}

\textbf{Context:} Creo que el contexto del proyecto queda bien explicado. Se describe el marco empresarial, se definen los conceptos clave del proyecto y se identifica con claridad el problema que se pretende resolver.

\textbf{Nota: Ben Assolit}

\textbf{Justificació:} He presentado las soluciones existentes en el mercado y he argumentado por qué ninguna de ellas se adapta plenamente a las necesidades del proyecto, justificando así el desarrollo de una solución propia.

\textbf{Nota: Ben Assolit}

\textbf{Abast:} Se han definido el objetivo general y los subobjetivos del sistema, junto con los requerimientos funcionales, no funcionales y las métricas de evaluación del código que el sistema utilizará para auditar la calidad del código generado.

\textbf{Nota: Ben Assolit}

\textbf{Metodologia i Rigor:} Se explica la metodología ágil adoptada (SCRUM) y cómo se adapta al desarrollo del proyecto, con sus sprints y reuniones de seguimiento. También se describen las herramientas utilizadas para el control del progreso.

\textbf{Nota: Ben Assolit}

\textbf{Referèncias:} Se han citado referencias variadas a lo largo del documento, combinando fuentes académicas —artículos de revista y conferencia sobre LLMs, agentes de código e ingeniería del software— con documentación técnica oficial de las herramientas y plataformas utilizadas. Muchas de las referencias, al ser a softwares o proyectos desarrollados por compañías, utilizan el nombre de la organización como autor en lugar de un nombre individual.

\textbf{Nota: Ben Assolit}

% ==================================================================
\section{Planificació Temporal}

\textbf{Descripció de les tasques:} Las tareas están desglosadas hasta el nivel de subtarea, con descripciones, estimaciones en horas, dependencias y recursos asignados. Se incluye una tabla resumen con toda esta información y un desglose de horas por rol.

\textbf{Nota: Ben Assolit}

\textbf{Estimacions i Gantt:} Las estimaciones se expresan íntegramente en horas y se distribuyen de forma coherente entre los cuatro roles definidos. El diagrama de Gantt refleja las dependencias entre tareas y la concurrencia de las actividades de documentación y seguimiento.

\textbf{Nota: Ben Assolit}

\textbf{Gestió del risc i plans alternatius:} Se han identificado seis riesgos específicos, cada uno con su plan alternativo, el impacto estimado en horas y los recursos adicionales necesarios. El coste económico de cada imprevisto está integrado en el presupuesto.

\textbf{Nota: Ben Assolit}

% ==================================================================
\section{Pressupost}

\textbf{Identificació dels costos:} Se han detallado todas las partidas del proyecto: costes de personal por actividad y rol, costes genéricos (infraestructura, software y hardware con amortización), contingencias e imprevistos vinculados a los riesgos identificados.

\textbf{Nota: Ben Assolit}

\textbf{Estimació dels costos:} Los costes están fundamentados en fuentes externas referenciadas. Se han aplicado correctamente los multiplicadores de Seguridad Social y diferenciado el coste por hora según el perfil. Considero que el presupuesto total resultante es realista para el alcance del proyecto.

\textbf{Nota: Ben Assolit}

\textbf{Control de gestió:} Se propone un sistema de control continuo alineado con la metodología ágil, revisando los indicadores al finalizar cada sprint. Se definen cuatro métricas de seguimiento con sus fórmulas y criterios de actuación según la magnitud de la desviación.

\textbf{Nota: Ben Assolit}

\textbf{Informe de Sostenibilitat:} Se incluye la autoevaluación de la competencia y el análisis de las tres dimensiones de sostenibilidad (económica, social y ambiental), cubriendo para cada una los apartados de PPP, Vida Útil y Riesgos.

\textbf{Nota: Ben Assolit}

% ==================================================================
\section{Format}

\textbf{Organització:} El documento presenta una estructura lógica con los apartados bien enlazados. La portada incluye todos los datos requeridos, y las tablas, figuras y ecuaciones están numeradas y referenciadas en el texto.

\textbf{Nota: Ben Assolit}

\textbf{Claredat:} Considero que el contenido se presenta de forma clara. El uso de listas, tablas comparativas, ecuaciones y el diagrama de Gantt facilita la comprensión de las distintas secciones.

\textbf{Nota: Ben Assolit}

\textbf{Redacció:} He revisado el documento en su totalidad y confío en que los errores ortográficos y gramaticales son mínimos.

\textbf{Nota: Ben Assolit}

\end{document}
